% 第9章第6节：一般矩阵奇异值的扰动（奇异值的误差区间，秩保证，到奇异矩阵的距离）。

\section{一般矩阵奇异值的扰动}
\label{sec:ls-singular-value-sensitivity}

秩、可逆性与最佳低秩逼近的误差都由奇异值决定。本节先把第~\ref{sec:ls-spectral-value-sensitivity}节的特征值区间移到奇异值上，再据此判断哪些关于秩的结论在误差预算下仍然可靠。

\subsection{奇异值的误差区间}

长方矩阵没有特征值，不能直接使用Weyl界。办法是把它放进一个对称矩阵，使后者的特征值恰好包含原矩阵的奇异值。对$\bm A\in\mathbb R^{m\times n}$，记$p=\min(m,n)$，奇异值为$\sigma_1(\bm A)\geq\cdots\geq\sigma_p(\bm A)\geq0$，定义$(m+n)$阶对称增广矩阵
\[
 \bm H(\bm A)=\begin{bmatrix}\bm0&\bm A\\\bm A^\top&\bm0\end{bmatrix}.
\]

\begin{lemma}[对称增广矩阵的谱]
\label{lem:ls-symmetric-augmentation-spectrum}
$\bm H(\bm A)$的特征值为$\pm\sigma_1(\bm A)$、$\ldots$、$\pm\sigma_p(\bm A)$，以及$|m-n|$个零。特别地，按降序排列时它的前$p$个特征值依次为$\sigma_1(\bm A),\ldots,\sigma_p(\bm A)$；并且$\lVert\bm H(\bm A)\rVert_2=\lVert\bm A\rVert_2$。
\end{lemma}
\begin{proof}
由定理~\ref{thm:ls-svd-existence}，取完整的SVD $\bm A=\bm U\bm\Sigma\bm V^\top$，列向量为$\bm u_1,\ldots,\bm u_m$与$\bm v_1,\ldots,\bm v_n$。对$i\leq p$，有$\bm A\bm v_i=\sigma_i\bm u_i$与$\bm A^\top\bm u_i=\bm V\bm\Sigma^\top\bm e_i=\sigma_i\bm v_i$。令
\[
 \bm w_i^{\pm}=\frac1{\sqrt2}\begin{bmatrix}\bm u_i\\\pm\bm v_i\end{bmatrix},
\]
用$(\bm a;\bm b)$表示把$\bm a$、$\bm b$上下拼接成的列向量，则
\[
 \begin{aligned}
 \bm H(\bm A)\bm w_i^{\pm}
 &=\tfrac1{\sqrt2}(\pm\bm A\bm v_i;\bm A^\top\bm u_i)\\
 &=\tfrac1{\sqrt2}(\pm\sigma_i\bm u_i;\sigma_i\bm v_i)=\pm\sigma_i\bm w_i^{\pm}.
 \end{aligned}
\]

再补出特征值$0$。若$m>n$，对$j>n$，$\bm\Sigma^\top$的第$j$列为零，故$\bm A^\top\bm u_j=\bm0$，向量$(\bm u_j;\bm0)$是特征值$0$的特征向量；若$n>m$，对$j>m$有$\bm A\bm v_j=\bm0$，向量$(\bm0;\bm v_j)$是特征值$0$的特征向量。

这$2p+|m-n|=m+n$个向量两两正交且长度为$1$。由$\bm u_i$、$\bm v_i$各自标准正交，$\bm w_i^{\pm}$与$\bm w_j^{\pm}$的内积为
\[
 \tfrac12\bigl(\bm u_i^\top\bm u_j\pm\bm v_i^\top\bm v_j\bigr),
\]
两者符号相同时取正号、相反时取负号。所以$i\neq j$时为零，$i=j$而符号相反时为$\tfrac12(1-1)=0$，符号相同时为$1$。补充的向量$(\bm u_j;\bm0)$或$(\bm0;\bm v_j)$只含$j>p$的分量，与所有$\bm w_i^{\pm}$及彼此都正交。

所以它们构成$\bm H(\bm A)$的标准正交特征基，列出了全部特征值（计重数）。$\sigma_i\geq0\geq-\sigma_j$，所以降序排列时前$p$个是$\sigma_1,\ldots,\sigma_p$。

最后，对$\bm z=(\bm y;\bm x)$，$\bm H(\bm A)\bm z=(\bm A\bm x;\bm A^\top\bm y)$。$\bm A^\top$与$\bm A$有相同的非零奇异值，所以$\lVert\bm A^\top\rVert_2=\lVert\bm A\rVert_2$，于是
\[
 \begin{aligned}
 \lVert\bm H(\bm A)\bm z\rVert_2^2
 &=\lVert\bm A\bm x\rVert_2^2+\lVert\bm A^\top\bm y\rVert_2^2\\
 &\leq\lVert\bm A\rVert_2^2\bigl(\lVert\bm x\rVert_2^2+\lVert\bm y\rVert_2^2\bigr)
 =\lVert\bm A\rVert_2^2\lVert\bm z\rVert_2^2 ,
 \end{aligned}
\]
而$\bm z=\bm w_1^+$时$\lVert\bm H(\bm A)\bm z\rVert_2=\sigma_1=\lVert\bm A\rVert_2$，故$\lVert\bm H(\bm A)\rVert_2=\lVert\bm A\rVert_2$。
\end{proof}

\begin{theorem}[奇异值的Weyl界]
\label{thm:ls-weyl-singular-value}
设$\bm A,\bm E\in\mathbb R^{m\times n}$，$p=\min(m,n)$。对$i=1,\ldots,p$，
\begin{equation}
 |\sigma_i(\bm A+\bm E)-\sigma_i(\bm A)|\leq\lVert\bm E\rVert_2 .
 \label{eq:matrix-analysis-singular-value-perturbation}
\end{equation}
\end{theorem}
\begin{proof}
增广是线性的：$\bm H(\bm A+\bm E)=\bm H(\bm A)+\bm H(\bm E)$，三者都实对称。对$\bm H(\bm A)$与扰动$\bm H(\bm E)$应用定理~\ref{thm:ls-weyl-eigenvalue}，第$i$个特征值的变化不超过$\lVert\bm H(\bm E)\rVert_2$。由引理~\ref{lem:ls-symmetric-augmentation-spectrum}，$i\leq p$时第$i$个特征值就是第$i$个奇异值，且$\lVert\bm H(\bm E)\rVert_2=\lVert\bm E\rVert_2$。
\end{proof}
奇异值非负，所以在误差预算$\lVert\bm E\rVert_2\leq\varepsilon$下，真实的第$i$个奇异值落在非负区间
\[
 \bigl[\max(0,\sigma_i(\bm A)-\varepsilon),\ \sigma_i(\bm A)+\varepsilon\bigr]
\]
内。

最佳低秩逼近的误差也随之受控。第~\ref{chap:matrix-approximation}章将证明（定理~\ref{thm:matrix-analysis-best-low-rank}），对$0\leq k<p$，秩不超过$k$的矩阵中，与$\bm A$的最小谱范数距离为$\sigma_{k+1}(\bm A)$。把定理~\ref{thm:ls-weyl-singular-value}用于$i=k+1$，扰动前后这个最优误差之差不超过$\varepsilon$。

要注意，这只说明最优误差的数值稳定；取得最优误差的低秩矩阵及其所用的奇异方向是否同样稳定，是另一个问题，留到第~\ref{sec:ls-singular-subspace-sensitivity}节。

\subsection{秩保证}

奇异值区间不含零，就能保证相应的奇异值为正，从而给出秩的下界。
\begin{theorem}[由奇异值保证秩]
\label{thm:ls-singular-value-rank-guarantee}
设$\bm A,\bm E\in\mathbb R^{m\times n}$，$\lVert\bm E\rVert_2\leq\varepsilon$，$1\leq k\leq\min(m,n)$。若$\sigma_k(\bm A)>\varepsilon$，则$\operatorname{rank}(\bm A+\bm E)\geq k$。特别地，$m\geq n$且$\sigma_n(\bm A)>\varepsilon$时，$\bm A+\bm E$满列秩。
\end{theorem}
\begin{proof}
由定理~\ref{thm:ls-weyl-singular-value}，
\[
 \sigma_k(\bm A+\bm E)\geq\sigma_k(\bm A)-\varepsilon>0 ,
\]
而奇异值降序排列，所以$\bm A+\bm E$至少有$k$个正奇异值。由定理~\ref{thm:ls-svd-existence}，正奇异值的个数等于秩。
\end{proof}
反方向的结论不成立：$\sigma_k(\bm A)\leq\varepsilon$只表示区间$[\max(0,\sigma_k(\bm A)-\varepsilon),\sigma_k(\bm A)+\varepsilon]$包含零，无法排除真实矩阵的第$k$个奇异值为零，但也不能由此断言真实矩阵秩亏。

这里要区分两种秩。精确秩对扰动不连续：$\operatorname{diag}(1,0)+t\bm e_2\bm e_2^\top$对每个$t\neq0$都可逆，秩为$2$，$t=0$时秩突然降为$1$，无论$|t|$多小。所以对带误差的矩阵，追问精确秩没有意义；可靠的是给定误差预算下的有效秩，即定理保证的下界“大于$\varepsilon$的奇异值个数”。

本节只给出由误差预算能够保证的结论；实际计算中怎样选取阈值、把大于阈值的奇异值个数作为数值秩，在第~\ref{sec:ls-pseudoinverse-rank-change}节结合求解任务讨论。

\subsection{到奇异矩阵的距离}
对可逆方阵，同样的想法回答：最小要多大的扰动，才能使它变成奇异矩阵？
\begin{theorem}[到奇异矩阵的距离]
\label{thm:ls-distance-to-singularity}
设$\bm A\in\mathbb R^{n\times n}$可逆，奇异值为$\sigma_1\geq\cdots\geq\sigma_n>0$，$\bm u_n$、$\bm v_n$是与$\sigma_n$对应的单位左、右奇异向量。则
\begin{equation}
 \min\{\lVert\bm E\rVert_2:\bm A+\bm E\text{ 奇异}\}=\sigma_n(\bm A),
 \label{eq:ls-distance-to-singularity}
\end{equation}
最小值在$\bm E_*=-\sigma_n\bm u_n\bm v_n^\top$处取到。相对于$\lVert\bm A\rVert_2$，这个距离为
\[
 \frac{\sigma_n(\bm A)}{\lVert\bm A\rVert_2}=\frac1{\kappa_2(\bm A)} .
\]
\end{theorem}
\begin{proof}
分两步证明：较小的扰动保持可逆，而$\bm E_*$恰好使矩阵奇异。

下界：设$\lVert\bm E\rVert_2<\sigma_n$。对任意$\bm x\neq\bm0$，由定理~\ref{thm:matrix-analysis-condition-number-bound}证明中的估计$\lVert\bm A\bm x\rVert_2\geq\sigma_n\lVert\bm x\rVert_2$，
\[
 \begin{aligned}
 \lVert(\bm A+\bm E)\bm x\rVert_2
 &\geq\lVert\bm A\bm x\rVert_2-\lVert\bm E\bm x\rVert_2\\
 &\geq(\sigma_n-\lVert\bm E\rVert_2)\lVert\bm x\rVert_2>0 ,
 \end{aligned}
\]
所以$\bm A+\bm E$的零空间只含零向量，它可逆。

可达：对任意$\bm x$，
\[
 \lVert\bm u_n\bm v_n^\top\bm x\rVert_2=|\bm v_n^\top\bm x|\leq\lVert\bm x\rVert_2 ,
\]
且$\bm x=\bm v_n$时取等号，故$\lVert\bm E_*\rVert_2=\sigma_n$。又
\[
 (\bm A+\bm E_*)\bm v_n=\sigma_n\bm u_n-\sigma_n\bm u_n=\bm0 ,
\]
所以$\bm A+\bm E_*$奇异。相对距离由$\lVert\bm A\rVert_2=\sigma_1$与$\kappa_2(\bm A)=\sigma_1/\sigma_n$得到。
\end{proof}
$n$阶方阵奇异，当且仅当它的秩不超过$n-1$，所以这个定理也是第~\ref{chap:matrix-approximation}章最佳低秩逼近在$k=n-1$时的特例；这里的直接证明只用到最小奇异值的下界。

这为第~\ref{sec:ls-perturbation}节的矩阵条件数提供了另一种解释：$\kappa_2(\bm A)$越大，按相对尺度计，$\bm A$离奇异矩阵越近。这个结论的条件是允许任意的谱范数扰动；若扰动必须保持某种结构，例如对称或稀疏，可选的扰动变少，本定理只给出距离的下界，相应的距离需要另行分析。

还要注意，扰动超过$\sigma_n$并不意味着每个这样的扰动都使矩阵奇异，例如$\bm E=\bm A$的范数为$\sigma_1\geq\sigma_n$，而$\bm A+\bm E=2\bm A$仍然可逆。

\begin{example}[对角矩阵的可逆性边界]
\label{ex:ls-diagonal-distance-to-singularity}
回到第~\ref{sec:ls-perturbation}节的$\bm A=\operatorname{diag}(1,s)$，$0<s<1$。它的奇异值为$1$与$s$，$\kappa_2(\bm A)=1/s$。由定理~\ref{thm:ls-distance-to-singularity}，误差预算$\varepsilon<s$时可以保证$\bm A+\bm E$可逆。

这个门槛不能再放宽：$\bm E=-s\bm e_2\bm e_2^\top$的范数恰为$s$，它把第二个对角元抵消掉，使$\bm A$变成奇异的$\operatorname{diag}(1,0)$。按相对尺度，这个距离是$s/1=1/\kappa_2(\bm A)$；$s$越小，条件数越大，$\bm A$离奇异矩阵也越近。
\end{example}
在这个例子中，$s$很小时矩阵离奇异很近，但它的第一奇异方向$\bm e_1$并不因此变得不稳定。第~\ref{sec:ls-singular-subspace-sensitivity}节将说明，保留第一奇异方向的稳定性由截断处的间隙$\sigma_1-\sigma_2=1-s$决定，而可逆性的边界由最小奇异值$s$决定，二者不能互相替代。

